\chapter{Overview}
\label{chap:overview}

% archon:covers MR4654610ComputingRiemannRochPolynomialsAndClassifyingHyperKahlerFourfolds.lean MR4654610ComputingRiemannRochPolynomialsAndClassifyingHyperKahlerFourfolds/Basic.lean

This chapter formalizes the main results of Debarre, Huybrechts, Macr\`i, and Voisin
\cite{MR4654610}, which proves O'Grady's conjecture: every hyper-K\"ahler fourfold of
\(\mathrm{K3}^{[2]}\) numerical type is of \(\mathrm{K3}^{[2]}\) deformation type.
The proof proceeds via the Beauville--Bogomolov--Fujiki form, the
Huybrechts--Riemann--Roch polynomial, topological constraints from Guan's classification,
and the SYZ conjecture in dimension~4.

%-------------------------------------------------------------------
\section{Introduction}
%-------------------------------------------------------------------

% SOURCE: [DHMV22], Introduction, p. 1 (read from references/MR4654610-hyperkahler-fourfolds.tex)
% SOURCE QUOTE: "A \hk\ manifold is a simply connected smooth compact K\"ahler manifold which carries
% a nowhere degenerate holomorphic 2-form (called a symplectic form) unique up to multiplication by a nonzero
% constant."
\begin{definition}[Hyper-K\"ahler fourfold]
  \label{def:hyp-kahler-fourfold}
  \lean{HK.HyperKaehlerFourfold}
  \textit{Source: Debarre--Huybrechts--Macr\`i--Voisin \cite{MR4654610}, Introduction.}
  A \emph{hyper-K\"ahler manifold} is a simply connected smooth compact K\"ahler manifold
  which carries a nowhere degenerate holomorphic 2-form (a symplectic form), unique up to
  multiplication by a nonzero scalar.  A \emph{hyper-K\"ahler fourfold} is a hyper-K\"ahler
  manifold of complex dimension \(4\).
\end{definition}

\begin{definition}[\(\mathrm{K3}^{[2]}\) numerical type]
  \label{def:k3-squared-numerical-type}
  \lean{HK.K3SquaredNumericalType}
  \uses{def:bbf-form, def:fujiki-constant}
  % SOURCE: [DHMV22], Introduction (read from references/MR4654610-hyperkahler-fourfolds.tex)
  % SOURCE QUOTE: "We say that a \hk\ fourfold $X$ is \emph{of $\mathrm{K3}^{[2]}$ numerical type} if, for some K3 surface~$S$, there exists an isomorphism of abelian groups $\psi\colon H^2(X,\Z)\isomto H^2(S^{[2]},\Z)$ such that, for all $\alpha\in H^2(X,\Z)$, we have $\int_X \alpha^4 = \int_{S^{[2]}} \psi(\alpha)^4$ (this is equivalent to asking that the lattices $(H^2(X,\Z),q_X) $ and $(H^2(S^{[2]},\Z),q_{S^{[2]}}) $ be isometric and the Fujiki constants be the same)."
  \textit{Source: Debarre--Huybrechts--Macr\`i--Voisin \cite{MR4654610}, Introduction.}
  A hyper-K\"ahler fourfold \(X\) is \emph{of \(\mathrm{K3}^{[2]}\) numerical type} if,
  for some K3 surface \(S\), there exists an isomorphism of abelian groups
  \(\psi \colon H^2(X, \mathbb{Z}) \to H^2(S^{[2]}, \mathbb{Z})\) such that
  \(\int_X \alpha^4 = \int_{S^{[2]}} \psi(\alpha)^4\) for all
  \(\alpha \in H^2(X, \mathbb{Z})\).  Equivalently, the BBF lattices and Fujiki constants
  of \(X\) and \(S^{[2]}\) coincide.
\end{definition}

\begin{definition}[\(\mathrm{K3}^{[2]}\) deformation type]
  \label{def:k3-squared-deformation-type}
  \lean{HK.K3SquaredDeformationType}
  \uses{def:hyp-kahler-fourfold}
  % SOURCE: [DHMV22], Introduction (read from references/MR4654610-hyperkahler-fourfolds.tex)
  % SOURCE QUOTE: "In \cite{og1}, O'Grady  conjectured  that a \hk\ fourfold   of  K3$^{[2]}$ numerical type is of~K3$^{[2]}$ deformation type..."
  \textit{Source: Debarre--Huybrechts--Macr\`i--Voisin \cite{MR4654610}, Introduction.}
  A hyper-K\"ahler fourfold \(X\) is \emph{of \(\mathrm{K3}^{[2]}\) deformation type} if
  \(X\) can be deformed as a compact complex manifold in a smooth family to the Hilbert
  scheme \(S^{[2]}\) of two points on some K3 surface \(S\).
\end{definition}

% SOURCE: [DHMV22], Introduction, Theorem (corresponding to Theorem~\ref{thm:RRpoly}) (read from references/MR4654610-hyperkahler-fourfolds.tex)
% SOURCE QUOTE: "\begin{theo}\label{thm:RRpoly}
% Let $X$ be a \hk\ fourfold with classes $\lll,\mm\in H^2(X,\Z)$ such that $\int_X \lll^4=0$ and $\int_X\lll^2\mm^2=2$.\
% The Huybrechts--Riemann--Roch polynomial of $X$ satisfies
% \[
% P_{RR,X}(2k)=\chi(\P^2,\cO_{\P^2}(k+1))=\binom{k+3}{2}.
% \]
% Furthermore, the Fujiki constant and the Hodge and Chern numbers of $X$ are those of the Hilbert square of a K3 surface.
% \end{theo}"
\begin{theorem}[HRR formula for fourfolds]
  \label{thm:hrr-formula-fourfold}
  \lean{HK.hrr_formula_fourfold}
  \uses{def:standard-triple, thm:conjecture14-fourfold, lem:binomial-coefficient}
  \textit{Source: Debarre--Huybrechts--Macr\`i--Voisin \cite{MR4654610}, Introduction (Theorem 1.5).}
  Let \(X\) be a hyper-K\"ahler fourfold with \(l, m \in H^2(X, \mathbb{Z})\) satisfying
  \(\int_X l^4 = 0\) and \(\int_X l^2 m^2 = 2\).  For all \(k \in \mathbb{Z}\),
  \[
    P_{RR,X}(2k) = \chi(\mathbb{P}^2, \mathcal{O}_{\mathbb{P}^2}(k+1)) = \binom{k+3}{2}.
  \]
\end{theorem}

\begin{proof}
  \uses{thm:conjecture14-fourfold, lem:binomial-coefficient}
  By \cref{thm:conjecture14-fourfold}, \(q_X(l) = 0\), \(q_X(l,m) = 1\), and
  \(P_{RR,X}(T) = \binom{T/2 + 3}{2}\).  Evaluate at \(T = 2k\):
  \(P_{RR,X}(2k) = \binom{k + 3}{2}\).
\end{proof}

% SOURCE: [DHMV22], Introduction, Theorem 1.1 (read from references/MR4654610-hyperkahler-fourfolds.tex)
% SOURCE QUOTE: "\begin{theo}\label{thm:MainThm}
% Let $X$ be a \hk\ fourfold.\
% Assume there are classes $\lll$, $\mm\in H^2(X,{\Z})$ such that $\int_X \lll^4=0$ and $\int_X\lll^2\mm^2=2$.\
% Then $X$ is of~$\mathrm{K3}^{[2]}$ deformation type.
% \end{theo}"
\begin{theorem}[Fourfolds with \(\int l^2 m^2 = 2\) are of \(\mathrm{K3}^{[2]}\) type]
  \label{thm:k3-deformation-type}
  \lean{HK.k3_deformation_type}
  \uses{thm:conjecture14-fourfold, prop:nef-case-sfz, prop:divisorial-contraction, thm:hrr-from-lagrangian}
  \textit{Source: Debarre--Huybrechts--Macr\`i--Voisin \cite{MR4654610}, Theorem 1.1.}
  Let \(X\) be a hyper-K\"ahler fourfold.  Assume there are classes
  \(l, m \in H^2(X, \mathbb{Z})\) with \(\int_X l^4 = 0\) and \(\int_X l^2 m^2 = 2\).
  Then \(X\) is of \(\mathrm{K3}^{[2]}\) deformation type.
\end{theorem}

% SOURCE QUOTE PROOF: "Let us explain  how we prove Theorem~\ref{thm:MainThm}.\
% The first step   is to show Conjecture~\ref{conjnum} in dimension $4$ (see Theorem~\ref{th43a}).
% ...
% A study of the effective cone of $X$ then shows that there are two cases.
% ...
% The second case, mostly studied in Section \ref{sec:DivisorialContraction}, is when $X$ admits a divisorial contraction and $M$ is not nef.\
%     In this case, we  first prove Proposition~\ref{prop:CaseC2} (a slightly weaker version of Theorem~\ref{thm:SYZ}), namely the existence of a Lagrangian fibration $f\colon X\to\P^2$ with $f^*{\cO_{\P^2}}(1)\isom L^{k_L}$, for some positive integer~$k_L$.\
%     We then use this weaker result to prove that $X$ is of  K3$^{[2]}$ deformation type (see Section~\ref{sec:OGrady})."
\begin{proof}
  \uses{thm:conjecture14-fourfold, prop:nef-case-sfz, prop:divisorial-contraction, thm:hrr-from-lagrangian}
  By \cref{thm:conjecture14-fourfold}, the Chern and Hodge numbers of \(X\) match those
  of \(\mathrm{K3}^{[2]}\).  By deformation theory \cite[Theorem~1.3]{Ma4}, one can
  deform to assume \(l, m\) are of type \((1,1)\) and \(\operatorname{NS}(X) = \mathbb{Z}l \oplus \mathbb{Z}m\)
  (very general case), forming a standard triple.  By \cref{prop:nef-case-sfz},
  case~(C1) does not occur.  In case~(C2), \cref{prop:divisorial-contraction} provides a
  Lagrangian fibration \(f \colon X \to \mathbb{P}^2\) with
  \(f^*\mathcal{O}_{\mathbb{P}^2}(1) \cong L^{k_L}\).  One shows \(k_L = 1\) using
  \cref{thm:hrr-from-lagrangian} and the matching of HRR polynomials, giving
  \(f^*\mathcal{O}_{\mathbb{P}^2}(1) \cong L\).  The base of this Lagrangian fibration
  identifies the deformation type with \(\mathrm{K3}^{[2]}\) by deformation invariance.
\end{proof}

% SOURCE: [DHMV22], Introduction, Corollary 1.2 (read from references/MR4654610-hyperkahler-fourfolds.tex)
% SOURCE QUOTE: "\begin{coro}[O'Grady's conjecture]\label{cor:OGrady}
% A \hk\ fourfold   of  $\mathrm{K3}^{[2]}$ numerical type  is of~$\mathrm{K3}^{[2]}$ deformation type.
% \end{coro}"
\begin{corollary}[O'Grady's conjecture]
  \label{cor:ogrady-conjecture}
  \lean{HK.ogrady_conjecture}
  \uses{def:k3-squared-numerical-type, thm:k3-deformation-type}
  \textit{Source: Debarre--Huybrechts--Macr\`i--Voisin \cite{MR4654610}, Corollary 1.2.}
  A hyper-K\"ahler fourfold of \(\mathrm{K3}^{[2]}\) numerical type is of
  \(\mathrm{K3}^{[2]}\) deformation type.
\end{corollary}

\begin{proof}
  \uses{def:k3-squared-numerical-type, thm:k3-deformation-type}
  By \cref{def:k3-squared-numerical-type}, the BBF lattice of \(X\) is isometric to that
  of \(S^{[2]}\).  In particular, there exist \(l, m \in H^2(X, \mathbb{Z})\) with
  \(\int_X l^4 = 0\) and \(\int_X l^2 m^2 = 2\).  Apply \cref{thm:k3-deformation-type}.
\end{proof}

% SOURCE: [DHMV22], Introduction, Theorem 1.3 (read from references/MR4654610-hyperkahler-fourfolds.tex)
% SOURCE QUOTE: "\begin{theo}\label{thm:SYZ}
% Let $X$ be a \hk\ fourfold and let $L$ be  a nef line bundle on $X$.\ Set
% $\lll\coloneqq c_1(L) \in H^2(X,\Z)$ and assume that there exists $\mm\in H^2(X,\Z)$ such that     $\int_X \lll^4=0$ and $\int_X\lll^2\mm^2=2$.\
% There exists a Lagrangian fibration $f\colon X\to\P^2$ with $f^*{\cO_{\P^2}}(1)\isom L$.
% \end{theo}"
\begin{theorem}[Lagrangian fibration (SYZ conjecture in dimension 4)]
  \label{thm:lagrangian-fibration}
  \lean{HK.lagrangian_fibration}
  \uses{thm:k3-deformation-type}
  \textit{Source: Debarre--Huybrechts--Macr\`i--Voisin \cite{MR4654610}, Theorem 1.3.}
  Let \(X\) be a hyper-K\"ahler fourfold and \(L\) a nef line bundle on \(X\).  Set
  \(l \coloneqq c_1(L)\) and assume there exists \(m \in H^2(X, \mathbb{Z})\) with
  \(\int_X l^4 = 0\) and \(\int_X l^2 m^2 = 2\).  There exists a Lagrangian fibration
  \(f \colon X \to \mathbb{P}^2\) with \(f^*\mathcal{O}_{\mathbb{P}^2}(1) \cong L\).
\end{theorem}

\begin{proof}
  \uses{thm:k3-deformation-type, thm:hrr-from-lagrangian}
  By \cref{thm:k3-deformation-type}, \(X\) is of \(\mathrm{K3}^{[2]}\) deformation type,
  so the SYZ conjecture holds for \(X\) \cite[Theorem~1.3]{Ma4}: since \(L\) is nef with
  \(c_1(L)^4 = 0\), a power of \(L\) is semi-ample.  The analysis of Section 8 of the
  paper shows that \(k_L = 1\), giving \(f^*\mathcal{O}_{\mathbb{P}^2}(1) \cong L\).
\end{proof}

%-------------------------------------------------------------------
\section{Review of hyper-K\"ahler manifolds}
%-------------------------------------------------------------------

\subsection{The Beauville--Bogomolov--Fujiki form and the Fujiki constant}

% SOURCE: [DHMV22], Section 2.1, p. 3 (read from references/MR4654610-hyperkahler-fourfolds.tex)
% SOURCE QUOTE: "There exists a canonical integral nondivisible quadratic form $q_X$ (the
% \emph{Beauville--Bogomolov--Fujiki form}) on $H^2(X,\Z)$ and a positive rational constant $c_X$ (the
% \emph{Fujiki constant}) such that
% \begin{equation}\label{fuj}
% \forall \alpha\in H^2(X,\R)\qquad \int_X \alpha^{2n}=c_X\, q_X(\alpha)^n
% \end{equation}
% (after extending $q_X$ to a quadratic form on $H^2(X,\R)$)."
\begin{definition}[Beauville--Bogomolov--Fujiki form]
  \label{def:bbf-form}
  \lean{HK.BBFForm}
  \uses{def:hyp-kahler-fourfold}
  \textit{Source: Debarre--Huybrechts--Macr\`i--Voisin \cite{MR4654610}, Section 2.1.}
  Let \(X\) be a hyper-K\"ahler manifold of dimension \(2n\).  The
  \emph{Beauville--Bogomolov--Fujiki form} is the canonical integral nondivisible quadratic
  form \(q_X\) on \(H^2(X, \mathbb{Z})\), extended to a quadratic form on
  \(H^2(X, \mathbb{R})\).  It satisfies \(q_X(h) > 0\) for all K\"ahler classes \(h\).
\end{definition}

\begin{definition}[Fujiki constant]
  \label{def:fujiki-constant}
  \lean{HK.fujikiConstant}
  \uses{def:bbf-form}
  \textit{Source: Debarre--Huybrechts--Macr\`i--Voisin \cite{MR4654610}, Section 2.1.}
  Let \(X\) be a hyper-K\"ahler manifold of dimension \(2n\).  The \emph{Fujiki constant}
  \(c_X\) is the unique positive rational number such that
  \(\int_X \alpha^{2n} = c_X\, q_X(\alpha)^n\) for all \(\alpha \in H^2(X, \mathbb{R})\).
\end{definition}

% SOURCE: [DHMV22], Section 2.1, equation (1) (read from references/MR4654610-hyperkahler-fourfolds.tex)
% SOURCE QUOTE: "\forall \alpha\in H^2(X,\R)\qquad \int_X \alpha^{2n}=c_X\, q_X(\alpha)^n"
\begin{lemma}[Fujiki relation]
  \label{lem:fujiki-relation}
  \lean{HK.fujiki_relation}
  \uses{def:bbf-form, def:fujiki-constant}
  \textit{Source: Debarre--Huybrechts--Macr\`i--Voisin \cite{MR4654610}, Section 2.1, equation (1).}
  Let \(X\) be a hyper-K\"ahler manifold of dimension \(2n\).  For all
  \(\alpha \in H^2(X, \mathbb{R})\),
  \[
    \int_X \alpha^{2n} = c_X\, q_X(\alpha)^n.
  \]
\end{lemma}

\begin{proof}
  \uses{def:bbf-form, def:fujiki-constant}
  This is the defining characterization of the Fujiki constant; see Huybrechts
  \cite[Proposition~23.14]{huy}.
\end{proof}

% SOURCE: [DHMV22], Section 2.1, p. 3, equation (2) (read from references/MR4654610-hyperkahler-fourfolds.tex)
% SOURCE QUOTE: "Assume $q_X(\alpha)=0$.\
% Then $\alpha^{n+1}=0$ (see for example \cite[Proposition~24.1]{huy}) and, by comparing the coefficients of $t^n$ in the relation
% \[
% \int_X (t\alpha+ \beta)^{2n}=c_Xq_X(t\alpha+ \beta)^n=c_X(2tq_X( \alpha, \beta)+q_X( \beta))^n,
% \]
% we get
% \begin{equation}\label{pol}
% \forall \beta\in H^2(X,\R)\qquad \frac{1}{2^n}\binom{2n}{n}\int_X \alpha^{n}\beta^{n}= c_Xq_X( \alpha, \beta)^n.
% \end{equation}"
\begin{lemma}[Polarized Fujiki relation]
  \label{lem:fujiki-polarized}
  \lean{HK.fujiki_polarized}
  \uses{lem:fujiki-relation}
  \textit{Source: Debarre--Huybrechts--Macr\`i--Voisin \cite{MR4654610}, Section 2.1, equation (2).}
  Let \(X\) be a hyper-K\"ahler manifold of dimension \(2n\).  Assume
  \(\alpha \in H^2(X, \mathbb{R})\) with \(q_X(\alpha) = 0\).  For all
  \(\beta \in H^2(X, \mathbb{R})\),
  \[
    \frac{1}{2^n}\binom{2n}{n}\int_X \alpha^n \beta^n = c_X\, q_X(\alpha, \beta)^n.
  \]
\end{lemma}

% SOURCE QUOTE PROOF: "Assume $q_X(\alpha)=0$.\
% Then $\alpha^{n+1}=0$ (see for example \cite[Proposition~24.1]{huy}) and, by comparing the coefficients of $t^n$ in the relation
% \[
% \int_X (t\alpha+ \beta)^{2n}=c_Xq_X(t\alpha+ \beta)^n=c_X(2tq_X( \alpha, \beta)+q_X( \beta))^n,
% \]
% we get
% \begin{equation}\label{pol}
% \forall \beta\in H^2(X,\R)\qquad \frac{1}{2^n}\binom{2n}{n}\int_X \alpha^{n}\beta^{n}= c_Xq_X( \alpha, \beta)^n.
% \end{equation}"
\begin{proof}
  \uses{lem:fujiki-relation}
  Since \(q_X(\alpha) = 0\), one has \(\alpha^{n+1} = 0\) by
  \cite[Proposition~24.1]{huy}.  Compare the coefficients of \(t^n\) in
  \(\int_X (t\alpha + \beta)^{2n} = c_X q_X(t\alpha + \beta)^n
    = c_X\bigl(2t\,q_X(\alpha, \beta) + q_X(\beta)\bigr)^n\).
\end{proof}

% SOURCE: [DHMV22], Section 2.1, equation (3) (read from references/MR4654610-hyperkahler-fourfolds.tex)
% SOURCE QUOTE: "This is a particular case of the polarization of the Fujiki relation~\eqref{fuj}, which in dimension 4 takes the form
%  \begin{equation}\label{eq:Fujiki3}
%  3\!\int_X \alpha_1\alpha_2\alpha_3\alpha_4 =c_X\bigl( q_X(\alpha_1,\alpha_2)q_X(\alpha_3,\alpha_4) + q_X(\alpha_1,\alpha_3)q_X(\alpha_2,\alpha_4) + q_X(\alpha_1,\alpha_4)q_X(\alpha_2,\alpha_3)\bigr)
% \end{equation}
% for all $\alpha_1, \alpha_2, \alpha_3,\alpha_4\in H^2(X,\Z)$."
\begin{lemma}[Fujiki relation in dimension 4]
  \label{lem:fujiki-fourfold}
  \lean{HK.fujiki_fourfold}
  \uses{lem:fujiki-polarized}
  \textit{Source: Debarre--Huybrechts--Macr\`i--Voisin \cite{MR4654610}, Section 2.1, equation (3).}
  Let \(X\) be a hyper-K\"ahler fourfold with Fujiki constant \(c_X\).  For all
  \(\alpha_1, \alpha_2, \alpha_3, \alpha_4 \in H^2(X, \mathbb{Z})\),
  \[
    3\int_X \alpha_1 \alpha_2 \alpha_3 \alpha_4
    = c_X\bigl(
      q_X(\alpha_1,\alpha_2)q_X(\alpha_3,\alpha_4)
      + q_X(\alpha_1,\alpha_3)q_X(\alpha_2,\alpha_4)
      + q_X(\alpha_1,\alpha_4)q_X(\alpha_2,\alpha_3)
    \bigr).
  \]
\end{lemma}

\begin{proof}
  \uses{lem:fujiki-polarized}
  Polarize the Fujiki relation at \(n = 2\).
\end{proof}

\subsection{The Huybrechts--Riemann--Roch polynomial}

% SOURCE: [DHMV22], Section 2.2, p. 4 (read from references/MR4654610-hyperkahler-fourfolds.tex)
% SOURCE QUOTE: "By \cite[Corollary~23.18]{huy}, there is a polynomial
% \begin{equation*}
% P_{RR,X}(T)=\sum_{i=0}^n  a_i T^i\in\Q[T]
% \end{equation*}
% of degree $n$  such that, for each line bundle $L$ on $X$, one has
% \begin{equation}\label{rr}
% \chi(X,L)=P_{RR,X}(q_X(c_1(L))).
% \end{equation}
% This polynomial has the following properties:
% \begin{enumerate}[{\rm (a)}]
% \item\label{ita} the  constant   term of $P_{RR,X}(T)$ is $\chi(X,\cO_X)=n+1$;
% \item\label{itb} the leading term of $P_{RR,X}(T)$ is $\frac{c_X}{(2n)!} T^n $;
% \item the coefficients of $P_{RR,X}(T)$ are all positive (\cite[Theorem~1.1]{jia}).
% \end{enumerate}"
\begin{definition}[Huybrechts--Riemann--Roch polynomial]
  \label{def:hrr-polynomial}
  \lean{HK.HRRPolynomial}
  \uses{def:bbf-form}
  \textit{Source: Debarre--Huybrechts--Macr\`i--Voisin \cite{MR4654610}, Section 2.2.}
  Let \(X\) be a hyper-K\"ahler manifold of dimension \(2n\).  The
  \emph{Huybrechts--Riemann--Roch polynomial} is the unique polynomial
  \(P_{RR,X}(T) = \sum_{i=0}^n a_i T^i \in \mathbb{Q}[T]\) of degree \(n\) such that
  \(\chi(X, L) = P_{RR,X}(q_X(c_1(L)))\) for every line bundle \(L\) on \(X\).  Its
  properties are: (a)~the constant term is \(\chi(X, \mathcal{O}_X) = n+1\); (b)~the
  leading term is \(\dfrac{c_X}{(2n)!}T^n\); (c)~all coefficients are positive.
\end{definition}

\begin{lemma}[Mathlib: binomial coefficient]
  \label{lem:binomial-coefficient}
  \lean{Nat.choose}
  \mathlibok
  \textit{Provided by Mathlib.}
  For nonneg\-ative integers \(n\) and \(k\), the binomial coefficient \(\binom{n+k}{k} \in \mathbb{N}\).
\end{lemma}

% NOTE: Exact Mathlib4 declaration name to be verified via LSP before Phase 2a dispatch.
% Likely candidate: Polynomial.eval_intCast_map (camelCase for cast components in Mathlib4).
\begin{lemma}[Mathlib: integer-valued polynomial at integers]
  \label{lem:integer-valued-polynomial}
  \lean{Polynomial.eval_intCast_map}
  \mathlibok
  \textit{Provided by Mathlib.}
  A polynomial over \(\mathbb{Z}\) evaluates to an integer at every integer argument.
\end{lemma}

% SOURCE: [DHMV22], Section 2.2, Lemma 2.1 (read from references/MR4654610-hyperkahler-fourfolds.tex)
% SOURCE QUOTE: "\begin{lemm}\label{lem11}
% Let $X$ be a \hk\ manifold.\ For every class $\alpha\in H^2(X,\Z)$, one has $P_{RR,X}(q_X(\alpha))\in\Z$.
% \end{lemm}"
\begin{lemma}[HRR polynomial is integer-valued on the lattice]
  \label{lem:hrr-integral}
  \lean{HK.hrr_integral}
  \uses{def:hrr-polynomial, lem:integer-valued-polynomial}
  \textit{Source: Debarre--Huybrechts--Macr\`i--Voisin \cite{MR4654610}, Lemma 2.1.}
  Let \(X\) be a hyper-K\"ahler manifold.  For every class \(\alpha \in H^2(X, \mathbb{Z})\),
  \[
    P_{RR,X}(q_X(\alpha)) \in \mathbb{Z}.
  \]
\end{lemma}

% SOURCE QUOTE PROOF: "\begin{proof}
% Since the period map is surjective (\cite[Proposition~25.12]{huy}), the class $\alpha$ becomes the first Chern class of some holomorphic line bundle on a deformation $X'$ of $X$, and~\eqref{rr} implies $P_{RR,X'}(q_{X'}(\alpha))\in \Z$.\ Since the polynomial~$P_{RR,X} $ and the form $q_X$ are deformation invariant, the lemma follows.
% \end{proof}"
\begin{proof}
  \uses{def:hrr-polynomial}
  By surjectivity of the period map \cite[Proposition~25.12]{huy}, the class \(\alpha\)
  deforms to \(c_1(L')\) for a holomorphic line bundle \(L'\) on a deformation \(X'\)
  of \(X\).  Then \(P_{RR,X'}(q_{X'}(\alpha)) = \chi(X', L') \in \mathbb{Z}\).  Since
  \(P_{RR,X}\) and \(q_X\) are deformation invariant, the claim follows.
\end{proof}

% SOURCE: [DHMV22], Section 2.2, Lemma 2.2 (read from references/MR4654610-hyperkahler-fourfolds.tex)
% SOURCE QUOTE: "\begin{lemm}\label{lemmaa}
% Let $X$ be a \hk\ manifold.\ Assume   that there is a  class~$\lll \in H^2(X,\Z)$ such that $\int_X \lll^{2n}=0$.\ For every $ \mm\in H^2(X,\Z)$, the number
% \begin{equation}\label{defa2}
% a\coloneqq \frac{1}{n!}\int_X \lll^{n}\mm^{n}
% \end{equation}
% is  an integer.
% \end{lemm}"
\begin{lemma}[The number \(a\) is an integer]
  \label{lem:a-is-integer}
  \lean{HK.a_is_integer}
  \uses{lem:fujiki-relation, lem:hrr-integral}
  \textit{Source: Debarre--Huybrechts--Macr\`i--Voisin \cite{MR4654610}, Lemma 2.2.}
  Let \(X\) be a hyper-K\"ahler manifold of dimension \(2n\).  Assume
  \(l \in H^2(X, \mathbb{Z})\) with \(\int_X l^{2n} = 0\).  For every
  \(m \in H^2(X, \mathbb{Z})\), the number
  \[
    a \coloneqq \frac{1}{n!}\int_X l^n m^n
  \]
  is an integer.
\end{lemma}

% SOURCE QUOTE PROOF: "\begin{proof}
% By~\eqref{fuj}, we have $q_X(\lll)=0$.\ Set
% \begin{equation}\label{defP}
% \forall k\in \Z\qquad P(k)\coloneqq P_{RR,X}(q_X(k\lll+ \mm))=P_{RR,X}(2kq_X(\lll,\mm)+q_X(\mm)) .
% \end{equation}
% Then $P$ is a  polynomial of degree $n$ whose leading coefficient is  (use~\eqref{pol} and property~\ref{itb} above)
% \begin{equation}\label{defia}
% \frac{c_X}{(2n)!} (2 q_X(\lll,\mm))^n=\frac{1}{(n!)^2}\int_X \lll^{n}\mm^{n}=\frac{a}{n!}.
% \end{equation}
% By Lemma~\ref{lem11}, the polynomial $P$ takes integral values on integers, hence $a$ is  an integer.
% \end{proof}"
\begin{proof}
  \uses{lem:fujiki-relation, lem:hrr-integral}
  The Fujiki relation gives \(q_X(l) = 0\).  Set
  \(P(k) \coloneqq P_{RR,X}(2k\,q_X(l,m) + q_X(m))\); this is a degree-\(n\) polynomial
  in \(k\) with leading coefficient \(a/n!\).
  By \cref{lem:hrr-integral}, \(P\) takes integer values at all integers, hence \(a\) is
  an integer.
\end{proof}

%-------------------------------------------------------------------
\section{Lagrangian fibrations}
%-------------------------------------------------------------------

\subsection{The Huybrechts--Riemann--Roch polynomial}

% SOURCE: [DHMV22], Section 3.1, Theorem 3.1 (read from references/MR4654610-hyperkahler-fourfolds.tex)
% SOURCE QUOTE: "\begin{theo}\label{th21}
% Let $X$ be a \hk\ manifold of dimension~$2n$ with a Lagrangian fibration $f\colon X\to \P^n$.\
% Set $\lll\coloneqq c_1(f^*\cO_{\P^n}(1))\in H^2(X,\Z)$ and assume that there exists $\mm\in H^2(X,\Z)$ such that $\int_X \lll^{n}\mm^{n}=n!$.\
% Then, $q_X(\lll,\mm)=\pm 1$, the quadratic form~$q_X$ is even, $c_X=(2n-1)!!$,
% \begin{equation*}
% P_{RR,X}(T)= \binom{\frac{T}{2}+n+1}{n},
% \end{equation*}
% and the sublattice $\Z\lll\oplus \Z\mm$ of $(H^2(X,\Z),q_X)$ is isomorphic to a hyperbolic plane.
% \end{theo}"
\begin{theorem}[HRR polynomial from a Lagrangian fibration with \(a=1\)]
  \label{thm:hrr-from-lagrangian}
  \lean{HK.hrr_from_lagrangian}
  \uses{def:bbf-form, def:hrr-polynomial, lem:polynomial-parity}
  \textit{Source: Debarre--Huybrechts--Macr\`i--Voisin \cite{MR4654610}, Theorem 3.1.}
  Let \(X\) be a hyper-K\"ahler manifold of dimension \(2n\) with a Lagrangian fibration
  \(f \colon X \to \mathbb{P}^n\).  Set
  \(l \coloneqq c_1(f^*\mathcal{O}_{\mathbb{P}^n}(1)) \in H^2(X, \mathbb{Z})\) and assume
  there exists \(m \in H^2(X, \mathbb{Z})\) with \(\int_X l^n m^n = n!\).  Then
  \(q_X(l,m) = \pm 1\), the quadratic form \(q_X\) is even, \(c_X = (2n-1)!!\),
  \[
    P_{RR,X}(T) = \binom{\tfrac{T}{2} + n + 1}{n},
  \]
  and the sublattice \(\mathbb{Z}l \oplus \mathbb{Z}m\) of \((H^2(X, \mathbb{Z}), q_X)\)
  is isomorphic to a hyperbolic plane.
\end{theorem}

\begin{proof}
  \uses{def:hrr-polynomial, lem:hrr-integral, lem:polynomial-parity}
  Replace \((X, L)\) by a deformation where \(\operatorname{NS}(X') = \mathbb{Z}l \oplus \mathbb{Z}m\)
  and \(L'\) is nef; a Lagrangian fibration \(f'\) on \(X'\) still exists by
  \cite[Theorem~1.2, Claims~3.1--3.2]{mat:isotropic}.  Make \(M\) ample by replacing it
  with \(L^k \otimes M\) for \(k \gg 0\).  Then \(f'_*M\) is locally free of rank \(a = 1\)
  on \(\mathbb{P}^n\), so \(f'_*M \cong \mathcal{O}_{\mathbb{P}^n}(d)\) for some \(d\)
  \cite[Section~3]{ort}.  The projection formula and definition of \(P_{RR,X}\) give
  \(P_{RR,X}(T) = \binom{d + (T - q_X(m))/(2q_X(l,m)) + n}{n}\).  Using
  \(P_{RR,X}(0) = n+1\) and positivity of coefficients forces \(q_X(m) = 0\) and \(d = 1\).
  Apply \cref{lem:polynomial-parity} with \(c = n!\), \(c' = 1\), \(q = 2q_X(l,m)\) to
  conclude \(q_X(l,m) = 1\).
\end{proof}

% SOURCE: [DHMV22], Section 3.1, Lemma 3.2 (read from references/MR4654610-hyperkahler-fourfolds.tex)
% SOURCE QUOTE: "\begin{lemm}\label{star}
% Let $X$ be a \hk\ manifold of dimension $2n$.\ Assume that there are positive integers $c$, $c'$, and $q$  such that $c'$ is divisible by no nontrivial $n$th powers and the polynomial $P(T)\coloneqq \frac{c}{c'}P_{RR,X}(qT)$ is  monic  with integral coefficients.\ Then, either $q=1$, or $q=2$  and the quadratic form $q_X$ is even.
% \end{lemm}"
\begin{lemma}[Polynomial parity]
  \label{lem:polynomial-parity}
  \lean{HK.polynomial_parity}
  \uses{lem:hrr-integral}
  \textit{Source: Debarre--Huybrechts--Macr\`i--Voisin \cite{MR4654610}, Lemma 3.2.}
  Let \(X\) be a hyper-K\"ahler manifold of dimension \(2n\).  Assume there are positive
  integers \(c\), \(c'\), and \(q\) such that \(c'\) is divisible by no nontrivial
  \(n\)-th powers and the polynomial \(P(T) \coloneqq \tfrac{c}{c'} P_{RR,X}(qT)\) is
  monic with integral coefficients.  Then either \(q = 1\), or \(q = 2\) and \(q_X\) is
  even.
\end{lemma}

% SOURCE QUOTE PROOF: "\begin{proof}
% Let $\alpha\in H^2(X,\Z)$ and write $\frac{ q_X(\alpha)}{q}=:\frac{r}{s}$, where $r$ and $s$ are relatively prime integers.\ One has
% \begin{equation*}
%     P_{RR,X}( q_X(\alpha))=\frac{c'}c\, P\Bigl(\frac{q_X(\alpha)}{q}\Bigr)
%     =\frac{c'}c\, P\Bigl(\frac{r}{s}\Bigr)
% \end{equation*}
% and this is an integer by Lemma~\ref{lem11}.\ Since $P(T)$ is monic  with integral coefficients, $s^nP(\frac{r}{s})$ is an integer congruent to $r^n$ modulo $s$, hence prime to $s^n$.\ But $c's^nP(\frac{r}{s})= cs^nP_{RR,X}( q_X(\alpha))$ is divisible by $s^n$, hence so is $c'$.\ Our hypothesis implies $s=\pm1$, which proves that  $q$ divides all values that $q_X$ takes on $H^2(X,\Z)$.\
% Since the  integral bilinear form associated with $q_X$ is  not  divisible, either $q=1$, or $q=2$  and the quadratic form $q_X$ is even.
% \end{proof}"
\begin{proof}
  \uses{lem:hrr-integral}
  For \(\alpha \in H^2(X, \mathbb{Z})\), write \(q_X(\alpha)/q = r/s\) in lowest terms.
  Since \(P\) is monic with integral coefficients, \(s^n P(r/s) \equiv r^n \pmod{s}\),
  so \(s^n P(r/s)\) is coprime to \(s^n\).  But
  \(c'\, s^n P(r/s) = c\, s^n P_{RR,X}(q_X(\alpha))\) is divisible by \(s^n\) (using
  \cref{lem:hrr-integral}), so \(s^n \mid c'\).  The hypothesis on \(c'\) forces
  \(s = \pm 1\), i.e.\ \(q \mid q_X(\alpha)\) for all \(\alpha \in H^2(X, \mathbb{Z})\).
  Since the associated bilinear form of \(q_X\) is not divisible, this gives \(q = 1\)
  or \(q = 2\) with \(q_X\) even.
\end{proof}

\subsection{The exceptional divisor}

% SOURCE: [DHMV22], Section 3.2, before Lemma 3.3 (read from references/MR4654610-hyperkahler-fourfolds.tex)
% SOURCE QUOTE: "Since $q_X( [E])<0$, the prime divisor $E$ is therefore exceptional in the sense of \cite[d\'ef.~3.10]{bou} and
% \cite[Definition~3.2]{mar} hence it spans an extremal ray in the effective cone $\Eff(X)$.\ Moreover, by~\cite[prop.~1.4 and rem.~4.3]{dru}, there is a birational isomorphism..."
% and "Moreover, $q_X([E])=-2$ and $q_X([E],\mm)=-1$; in particular, $M$ is not nef."
\begin{definition}[Exceptional divisor of a standard triple]
  \label{def:exceptional-divisor}
  \lean{HK.exceptionalDivisor}
  \uses{def:standard-triple, def:bbf-form}
  \textit{Source: Debarre--Huybrechts--Macr\`i--Voisin \cite{MR4654610}, Section 3.2.}
  Let \((X, l, m)\) be a standard triple with
  \(\operatorname{NS}(X) = \mathbb{Z}l \oplus \mathbb{Z}m\).  The
  \emph{exceptional divisor} \(E\) is the unique prime divisor in \(|L^{-1} \otimes M|\)
  satisfying \(q_X([E]) = -2\) and \(q_X([E], m) = -1\).
\end{definition}

%-------------------------------------------------------------------
\section{Conjecture 1.4 for hyper-K\"ahler fourfolds}
%-------------------------------------------------------------------

% SOURCE: [DHMV22], Section 4, equation (5) (read from references/MR4654610-hyperkahler-fourfolds.tex)
% SOURCE QUOTE: "Following \cite[Section~2.4]{jia}, we
%  set
% \begin{equation}\label{ax}
% A_X\coloneqq\int_X\td^{1/2}(X)= \frac 1 {5760} ( 7c_2^2(X) -4c_4(X) ) = \frac 18 \Bigl(7-\frac 1 {432} c_4(X)\Bigr).
% \end{equation}"
\begin{definition}[Topological constant \(A_X\)]
  \label{def:topological-constant}
  \lean{HK.topologicalConstant}
  \uses{def:hyp-kahler-fourfold}
  \textit{Source: Debarre--Huybrechts--Macr\`i--Voisin \cite{MR4654610}, Section 4, equation (5).}
  For a hyper-K\"ahler fourfold \(X\), define the topological constant
  \[
    A_X \coloneqq \int_X \mathrm{td}^{1/2}(X)
    = \frac{1}{5760}\bigl(7c_2(X)^2 - 4c_4(X)\bigr).
  \]
\end{definition}

% SOURCE: [DHMV22], Section 4, Lemma 4.1 (read from references/MR4654610-hyperkahler-fourfolds.tex)
% SOURCE QUOTE: "\begin{lemm}\label{lemmguan}
% Let $X$ be a \hk\ fourfold.\ Then,
% \begin{enumerate}[{\rm (a)}]
% \item\label{aax} either $b_2(X) =23$,  $b_3(X) =0$, and the Hodge numbers of $X$ are those of the Hilbert square of a K3 surface, in which case $ c_4(X)= 324$ and $A_X=\frac{25}{32}$,
% \item\label{bax} or $b_2(X) \le 8$, in which case $c_4(X)\le 144$ and  $\frac{5}{6}\le A_X\le \frac{131}{144}$.
%  \end{enumerate}
% In particular, if $t\in [0,\frac13)$, then $4A_X-t $ is an integer only when $t=\frac18$ and $A_X=\frac{25}{32}$.
% \end{lemm}"
\begin{lemma}[Betti numbers of hyper-K\"ahler fourfolds]
  \label{lem:betti-numbers}
  \lean{HK.betti_numbers}
  \uses{def:topological-constant}
  \textit{Source: Debarre--Huybrechts--Macr\`i--Voisin \cite{MR4654610}, Lemma 4.1.}
  Let \(X\) be a hyper-K\"ahler fourfold.  Then either:
  \begin{enumerate}
    \item[(a)] \(b_2(X) = 23\), \(b_3(X) = 0\), and the Hodge numbers of \(X\) are those
      of the Hilbert square of a K3 surface, with \(c_4(X) = 324\) and
      \(A_X = 25/32\); or
    \item[(b)] \(b_2(X) \le 8\), \(c_4(X) \le 144\), and \(5/6 \le A_X \le 131/144\).
  \end{enumerate}
  In particular, if \(t \in [0, 1/3)\), then \(4A_X - t\) is an integer only when
  \(t = 1/8\) and \(A_X = 25/32\).
\end{lemma}

\begin{proof}
  \uses{def:topological-constant}
  Apply the relation \(c_4(X) = 3(4b_2(X) + 16 - b_3(X))\) and the classification of
  possible Betti numbers for hyper-K\"ahler fourfolds from Guan
  \cite[Main Theorem]{gua}.
\end{proof}

% SOURCE: [DHMV22], Section 4, Lemma 4.2 (read from references/MR4654610-hyperkahler-fourfolds.tex)
% SOURCE QUOTE: "\begin{lemm}\label{lemmax}
% The number $\sqrt{2aA_X} $ is   rational.
% \end{lemm}"
\begin{lemma}[Square-root rationality]
  \label{lem:sqrt-rationality}
  \lean{HK.sqrt_rationality}
  \uses{def:topological-constant, lem:a-is-integer, def:hrr-polynomial}
  \textit{Source: Debarre--Huybrechts--Macr\`i--Voisin \cite{MR4654610}, Lemma 4.2.}
  Let \(X\) be a hyper-K\"ahler fourfold with \(l, m \in H^2(X, \mathbb{Z})\) satisfying
  \(q_X(l) = 0\) and \(a \coloneqq \tfrac{1}{2}\int_X l^2 m^2\).  Then
  \(\sqrt{2aA_X} \in \mathbb{Q}\).
\end{lemma}

% SOURCE QUOTE PROOF: "\begin{proof}
% By~\eqref{defia}, we have $c_Xq_X(\lll,\mm)^2=3a$.\ Using Section~\ref{subsec:RRpolynomial} and~\cite[Lemma~5.7]{jia}, we obtain
% \begin{equation}\label{prr}
%    P_{RR,X}(T)=\frac{c_X}{24}T^2+ T\sqrt{\frac23c_XA_X} +3=\frac{a}{8}\Bigl(\frac{T}{q_X(\lll,\mm)}\Bigr)^2+ \sqrt{2aA_X}\Bigl(\frac{T}{q_X(\lll,\mm)}\Bigr) +3.
% \end{equation}
% This implies that $\sqrt{2aA_X} $ is a rational number.
% \end{proof}"
\begin{proof}
  \uses{def:hrr-polynomial, lem:a-is-integer}
  From the Fujiki relation and \cref{lem:a-is-integer}, one has \(c_X q_X(l,m)^2 = 3a\).
  Using \cite[Lemma~5.7]{jia},
  \[
    P_{RR,X}(T)
    = \frac{c_X}{24}T^2 + T\sqrt{\tfrac{2}{3}c_X A_X} + 3
    = \frac{a}{8}\Bigl(\frac{T}{q_X(l,m)}\Bigr)^{\!2}
      + \sqrt{2aA_X}\,\frac{T}{q_X(l,m)} + 3.
  \]
  Since \(P_{RR,X}\) has rational coefficients, \(\sqrt{2aA_X}\) is rational.
\end{proof}

% SOURCE: [DHMV22], Section 4, Theorem 4.3 (read from references/MR4654610-hyperkahler-fourfolds.tex)
% SOURCE QUOTE: "\begin{theo}\label{th43a}
% Let $X$ be a  \hk\ fourfold with classes~$\lll,\mm\in H^2(X,\Z)$ such that $\int_X \lll^{4}=0$
% and   $\int_X \lll^{2}\mm^{2}=2$.\
%  The Chern and Hodge numbers of $X$ are those of the Hilbert square of a K3 surface  and all the  conclusions of Theorem~\ref{th21} hold.
% \end{theo}"
\begin{theorem}[Conjecture 1.4 in dimension 4]
  \label{thm:conjecture14-fourfold}
  \lean{HK.conjecture14_fourfold}
  \uses{lem:betti-numbers, lem:sqrt-rationality, lem:polynomial-parity, def:hrr-polynomial}
  \textit{Source: Debarre--Huybrechts--Macr\`i--Voisin \cite{MR4654610}, Theorem 4.3.}
  Let \(X\) be a hyper-K\"ahler fourfold with \(l, m \in H^2(X, \mathbb{Z})\) satisfying
  \(\int_X l^4 = 0\) and \(\int_X l^2 m^2 = 2\).  Then the Chern and Hodge numbers of
  \(X\) are those of the Hilbert square of a K3 surface, and \(q_X(l,m) = 1\), \(q_X\) is
  even, \(c_X = 3\), and
  \[
    P_{RR,X}(T) = \binom{\tfrac{T}{2} + 3}{2}.
  \]
\end{theorem}

% SOURCE QUOTE PROOF: "\begin{proof}
% Changing   $\mm$ into $\pm\mm+r\lll$, for some $r\in\Z$, if necessary, we may assume $-q_X(\lll,\mm)< q_X(\mm)\le q_X(\lll,\mm)$ or, equivalently, $\gamma\coloneqq \frac{q_X(\mm)}{q_X(\lll,\mm)}\in (-1,1] $.\
% As in~\eqref{defP}, we introduce the polynomial
% $$P(k)\coloneqq P_{RR,X}(q_X(k\lll+ \mm))=P_{RR,X}(2kq_X(\lll,\mm)+q_X(\mm)).$$
% Using~\eqref{prr}, we compute
% \begin{align*}
%  P(k)&= \frac{1}{8}(2k+\gamma)^2+ \sqrt{2A_X}(2k+\gamma) +3  \\
%  &=\frac{1}{2}k^2+\Bigl( \frac{1}{2}\gamma+2\sqrt{2A_X} \Bigr)k+\frac{1}{8}\gamma^2+ \gamma\sqrt{2A_X}+3\\
%  &=:\frac{1}{2}k^2 +  b  k +c.
% \end{align*}
% Since $P$ takes integral values on integers, $\frac{1}2+b=P(1)-P(0)$ and $c=P(0)$ are integers; we write $b=\frac12+b'$, with $b'\in\Z$.
% We also note that
% \begin{align*}
%     4A_X-\frac{b^2}{2}=3-c
% \end{align*}
% is an integer, hence so is $4A_X-\frac{1}{8}$.\ By Lemma~\ref{lemmguan}, this is only possible when $A_X=\frac{25}{32}$ and the Chern and Hodge numbers of $X$ are those of the Hilbert square of a K3 surface.\
% So $\gamma=0$,  hence $q_X(\mm)=0$, $b=\frac52$, and $c=3$.\
% By Lemma~\ref{star} (applied with $c=2$, $c'=1$, and $q=2q_X(\lll,\mm)$), we get $q_X(\lll,\mm)=1$, the quadratic form $q_X$ is even, and all the conclusions of Theorem~\ref{th21} hold.
% \end{proof}"
\begin{proof}
  \uses{lem:betti-numbers, lem:sqrt-rationality, lem:polynomial-parity}
  Replacing \(m\) by \(\pm m + rl\) for \(r \in \mathbb{Z}\), assume
  \(\gamma \coloneqq q_X(m)/q_X(l,m) \in (-1, 1]\).  Set
  \(P(k) \coloneqq P_{RR,X}(2k\,q_X(l,m) + q_X(m))\); using \cref{lem:sqrt-rationality}
  with \(a = 1\),
  \[
    P(k) = \tfrac{1}{2}k^2 + bk + c,
    \quad
    b = \tfrac{1}{2}\gamma + 2\sqrt{2A_X},
    \quad
    c = \tfrac{1}{8}\gamma^2 + \gamma\sqrt{2A_X} + 3.
  \]
  Since \(P\) takes integer values at integers, \(4A_X - b^2/2 = 3 - c\) is an integer,
  so \(4A_X - 1/8 \in \mathbb{Z}\).  By \cref{lem:betti-numbers}, this forces
  \(A_X = 25/32\), giving the K3\(^{[2]}\) Hodge/Chern numbers.  Then \(\gamma = 0\),
  \(q_X(m) = 0\), \(c = 3\), and \(P_{RR,X}(2k\,q_X(l,m)) = \binom{k+3}{2}\).
  Apply \cref{lem:polynomial-parity} with parameters \(c = 2\), \(c' = 1\),
  \(q = 2q_X(l,m)\) to conclude \(q_X(l,m) = 1\) and \(q_X\) even.
\end{proof}

%-------------------------------------------------------------------
\section{On the SYZ conjecture in dimension 4}
%-------------------------------------------------------------------

\subsection{Cones of divisors}

% SOURCE: [DHMV22], Section 5.2, Lemma 5.2 (read from references/MR4654610-hyperkahler-fourfolds.tex)
% SOURCE QUOTE: "\begin{lemm}\label{lem:PrimeExceptional}
% Let $E$ be a prime exceptional divisor on $X$.\
% We have $[E]=\pm(-\lll+\mm)$ in $\NS(X)$.
% \end{lemm}"
\begin{lemma}[Class of a prime exceptional divisor]
  \label{lem:prime-exceptional-class}
  \lean{HK.prime_exceptional_class}
  \uses{def:exceptional-divisor}
  \textit{Source: Debarre--Huybrechts--Macr\`i--Voisin \cite{MR4654610}, Lemma 5.2.}
  Let \((X, l, m)\) be a standard triple with
  \(\operatorname{NS}(X) = \mathbb{Z}l \oplus \mathbb{Z}m\), and let \(E\) be a prime
  exceptional divisor on \(X\).  Then \([E] = \pm(-l + m)\) in \(\operatorname{NS}(X)\).
\end{lemma}

% SOURCE QUOTE PROOF: "\begin{proof}
% Let us write $[E]=t\lll+u\mm$, with $t,u\in\Z$.
% By~\cite[Corollary~3.6]{mar}, the class
% \[
% -2\frac{q_X([E],\bullet)}{q_X([E])} = -\frac{q_X(t\lll+u\mm,\bullet)}{tuq_X(\lll,\mm)}
% \]
% is
% in $H_2(X,\Z)$.\
% By applying it to $\lll$ and $\mm$, we get $|t|=|u|=1$.\ Since $q_X([E])<0$, we have $t=-u$, as we wanted.
% \end{proof}"
\begin{proof}
  \uses{def:exceptional-divisor}
  Write \([E] = tl + um\) with \(t, u \in \mathbb{Z}\).  By
  \cite[Corollary~3.6]{mar}, the class \(-2q_X([E], \cdot)/q_X([E])\) lies in
  \(H_2(X, \mathbb{Z})\).  Applying this to \(l\) and \(m\) separately gives
  \(|t| = |u| = 1\).  Since \(q_X([E]) = 2tu\,q_X(l,m) = 2tu < 0\), we have
  \(tu < 0\); hence \(t = -u\), giving \([E] = \pm(-l + m)\).
\end{proof}

\subsection{The two key propositions}

% SOURCE: [DHMV22], Section 5.2, equations (16) and (16') (read from references/MR4654610-hyperkahler-fourfolds.tex)
% SOURCE QUOTE: "In the rest of this article, we will be mostly concerned with triples $(X,\lll,\mm)$ satisfying the following set of properties:
% \begin{equation}\label{hypo}
%     \left\{
%     \begin{array}{l}
%     X \hbox{ is a  \hk\ fourfold;}\\
%    \lll, \mm\in \NS(X);\\
%    \int_X \lll^{4}=0 \hbox { and }\int_X \lll^{2}\mm^{2}=2.
%     \end{array}
%     \right.
% \end{equation}"
\begin{definition}[Standard triple]
  \label{def:standard-triple}
  \lean{HK.StandardTriple}
  \uses{thm:conjecture14-fourfold}
  \textit{Source: Debarre--Huybrechts--Macr\`i--Voisin \cite{MR4654610}, Section 5.2.}
  A \emph{standard triple} \((X, l, m)\) consists of a hyper-K\"ahler fourfold \(X\) and
  classes \(l, m \in \operatorname{NS}(X)\) satisfying \(\int_X l^4 = 0\) and
  \(\int_X l^2 m^2 = 2\).  By \cref{thm:conjecture14-fourfold}, such a triple also
  satisfies (after possibly replacing \(m\) by \(-m\) plus a multiple of \(l\)):
  \(q_X(l) = q_X(m) = 0\), \(q_X(l,m) = 1\), \(P_{RR,X}(T) = \binom{T/2 + 3}{2}\),
  \(q_X\) is even, and \(c_X = 3\).
\end{definition}

% SOURCE: [DHMV22], Section 5.2, Proposition 5.5 (read from references/MR4654610-hyperkahler-fourfolds.tex)
% SOURCE QUOTE: "\begin{prop}\label{lem:NefCone4folds}
% Let $(X,\lll,\mm)$ be a triple satisfying~\eqref{hypo} and such that~$\NS(X)=\Z\lll\oplus\Z\mm$.\
% The nef and movable cones of $X$ coincide.\
% Equivalently, $X$ has no nontrivial \hk\ birational models.
% \end{prop}"
\begin{lemma}[Nef and movable cones coincide]
  \label{lem:nef-movable-coincide}
  \lean{HK.nef_movable_coincide}
  \uses{def:standard-triple}
  \textit{Source: Debarre--Huybrechts--Macr\`i--Voisin \cite{MR4654610}, Proposition 5.5.}
  Let \((X, l, m)\) be a very general standard triple with
  \(\operatorname{NS}(X) = \mathbb{Z}l \oplus \mathbb{Z}m\).  The nef and movable cones
  of \(X\) coincide; equivalently, \(X\) has no nontrivial hyper-K\"ahler birational
  models.
\end{lemma}

% SOURCE QUOTE PROOF: "\begin{proof}
% We follow the proof of \cite[Theorem~22]{hats}, whose argument we briefly review.\
% We assume for a contradiction that the nef and movable cones of $X$ are different.\
% By \cite[Theorem 1.1]{wiwi}, there exists a Mukai flop on $X$, hence a Lagrangian plane $\P^2\subset X$.\
% ...
% Solving~\eqref{eq:NefCone4folds1}, \eqref{eq:NefCone4folds2}, and \eqref{eq:NefCone4folds3} for $t$ and $u$, we get the equation
% \[
% 92 q_X(A)^2 + 20  q_X(A) - 525 = 0,
% \]
% which has no integral solutions.\
% This a contradiction, which proves the proposition.
% \end{proof}"
\begin{proof}
  \uses{def:standard-triple}
  Assume for contradiction that the nef and movable cones differ.  By
  \cite[Theorem~1.1]{wiwi}, there exists a Mukai flop, yielding a Lagrangian
  \(\mathbb{P}^2 \subset X\).  Choose \(A \in \operatorname{NS}(X)\) dual to the line
  class \(\ell \in H_2(X, \mathbb{Z})\) via \(q_X(A, B) = B \cdot \ell\).  Writing
  \([\mathbb{P}^2] = tA^2 + u q_X^\vee \in H^4(X, \mathbb{Q})\) and computing the
  intersection numbers \([\mathbb{P}^2]^2\), \(c_2(X) \cdot [\mathbb{P}^2]\),
  \(A^2 \cdot [\mathbb{P}^2]\) using \(c_X = 3\) and the formulas from
  \cref{def:standard-triple}, one arrives at the Diophantine equation
  \(92\,q_X(A)^2 + 20\,q_X(A) - 525 = 0\), which has no integral solutions.
\end{proof}

% SOURCE: [DHMV22], Section 5.2, Proposition 5.6 (read from references/MR4654610-hyperkahler-fourfolds.tex)
% SOURCE QUOTE: "\begin{prop}\label{prop:CaseC1}
% Let $(X,\lll,\mm)$ be a very general triple satisfying~\eqref{hypo}
% and assume in addition that we are in case~{\rm\ref{caseC1}}.\
%  Then one of the following  statements holds:
% \begin{enumerate}[{\rm (a)}]
%        \item  $H^0(X,L)\ne 0$, $H^0(X,M)\ne 0$, and  there is a Lagrangian fibration \mbox{$f\colon X\to\P^2$} with $f^*{\cO_{\P^2}}(1)\isom L$ or $M$;\label{prop64a}
%     \item  at least one of $H^0(X,L)$ or $H^0(X,M)$ is trivial and  the image of the rational map $\phi_{L\otimes M}\colon X\dra \P^5$ is rationally connected.\label{prop64b}
% \end{enumerate}
% \end{prop}"
\begin{proposition}[Nef case: structure of the linear systems]
  \label{prop:nef-case-sfz}
  \lean{HK.nef_case}
  \uses{def:standard-triple, lem:nef-movable-coincide, def:exceptional-divisor}
  \textit{Source: Debarre--Huybrechts--Macr\`i--Voisin \cite{MR4654610}, Proposition 5.6.}
  Let \((X, l, m)\) be a very general standard triple in case~(C1): \(m\) is nef.  Then
  one of the following holds:
  \begin{enumerate}
    \item[(a)] \(H^0(X, L) \ne 0\), \(H^0(X, M) \ne 0\), and there is a Lagrangian
      fibration \(f \colon X \to \mathbb{P}^2\) with
      \(f^*\mathcal{O}_{\mathbb{P}^2}(1) \cong L\) or \(M\);
    \item[(b)] at least one of \(H^0(X, L)\) or \(H^0(X, M)\) vanishes, and the image of
      the rational map \(\phi_{L \otimes M} \colon X \dashrightarrow \mathbb{P}^5\) is
      rationally connected.
  \end{enumerate}
  In particular, case~(C1) does not occur: case~(b) is excluded by
  \cite[Theorem~4.2]{voisinfib}, while case~(a) contradicts \cref{def:exceptional-divisor}
  (in a Lagrangian fibration both \(L\) and \(M\) cannot be nef simultaneously).
\end{proposition}

% SOURCE QUOTE PROOF: "6.3. Proof of Proposition 5.6. If H^0(X,L) and H^0(X,M) are both nonzero,
% by Lemma 6.1 and Lemma 5.1, after possibly permuting L and M, the line bundle L is globally
% generated and thus gives a Lagrangian fibration f: X -> P^2 with f*O_{P^2}(1) =~ L. So we
% are in case (a) of Proposition 5.6. If either H^0(X,L) or H^0(X,M) is zero, the image of the
% rational map phi_{L tensor M} is rationally connected by Proposition 6.3 and we are in case (b)."
\begin{proof}
  \uses{def:standard-triple, lem:nef-movable-coincide}
  \textit{Source: Debarre--Huybrechts--Macr\`i--Voisin \cite{MR4654610}, Section~6.3.}
  In case~(C1), both \(l\) and \(m\) are nef and isotropic with
  \(h^0(X, L^p \otimes M^q) = \binom{pq+3}{2}\) for \(p, q > 0\) by Kodaira vanishing.

  \textbf{Sub-case~(a):} If \(H^0(X,L) \ne 0\) and \(H^0(X,M) \ne 0\), then by the
  intersection computation (Lemma~6.1 of the paper) at least one of
  \(H^0(X,L)\) or \(H^0(X,M)\) has dimension \(\ge 2\).  By Lemma~5.1 (applied since
  \(\kappa(X,L) > 0\)) the line bundle \(L\) (or \(M\)) is globally generated and defines
  a Lagrangian fibration \(f \colon X \to \mathbb{P}^2\).  But then by the discussion
  of Section~3.2, both \(L\) and \(M\) cannot be nef simultaneously in a Lagrangian
  fibration (the exceptional divisor \(E\) of class \(-l+m\) must be prime and irreducible
  by \cref{def:exceptional-divisor}).  Contradiction.

  \textbf{Sub-case~(b):} If \(H^0(X,L) = 0\) or \(H^0(X,M) = 0\), the image of
  \(\phi_{L \otimes M} \colon X \dashrightarrow \mathbb{P}^5\) is rationally connected
  (Proposition~6.3 of the paper, proved via Castelnuovo's lemma and Lemma~6.6).
  This contradicts \cite[Theorem~4.2]{voisinfib} for very general \(X\).

  Hence case~(C1) does not occur.
\end{proof}

% SOURCE: [DHMV22], Section 5.2, Proposition 5.8 (read from references/MR4654610-hyperkahler-fourfolds.tex)
% SOURCE QUOTE: "\begin{prop}\label{prop:CaseC2}
% Let $(X,\lll,\mm)$ be a very general triple
% satisfying~\eqref{hypo} and assume in addition that we are in case~{\rm\ref{caseC2}}.\
%  There exist a positive integer~$k_L$ and a Lagrangian fibration \mbox{$f\colon X\to\P^2$} such that $f^*{\cO_{\P^2}}(1)\isom L^{\otp {k_L}}$.
% \end{prop}"
\begin{proposition}[Divisorial contraction case: Lagrangian fibration]
  \label{prop:divisorial-contraction}
  \lean{HK.divisorial_contraction}
  \uses{def:standard-triple, def:exceptional-divisor, lem:nef-movable-coincide}
  \textit{Source: Debarre--Huybrechts--Macr\`i--Voisin \cite{MR4654610}, Proposition 5.8.}
  Let \((X, l, m)\) be a very general standard triple in case~(C2): there exists a prime
  exceptional divisor \(E\) of class \(-l + m\).  There exist a positive integer
  \(k_L\) and a Lagrangian fibration \(f \colon X \to \mathbb{P}^2\) such that
  \(f^*\mathcal{O}_{\mathbb{P}^2}(1) \cong L^{k_L}\).
\end{proposition}

% SOURCE QUOTE PROOF: "Proof of Proposition 5.8. As observed in Section 5.1, we only have to
% show h^0(X,L^2) >= 2. This inequality follows from Lemma 7.3 by deformation and
% specialization: as we explained in Section 3.2, the reflection that permutes l and m is a
% monodromy operator. This means that one can deform the pair (X,L) into the pair (X,M)
% through a family (X, L) -> T. By semi-continuity, Lemma 7.3 implies h^2(X_t, L_t^2) =
% h^4(X_t, L_t^2) = 0 for t in T general. The Riemann-Roch argument of (40) then implies
% h^0(X_t, L_t^2) >= 3 for t in T general. By upper semi-continuity, h^0(X,L^2) >= 3."
\begin{proof}
  \uses{def:standard-triple, def:exceptional-divisor}
  \textit{Source: Debarre--Huybrechts--Macr\`i--Voisin \cite{MR4654610}, Section~7.}
  In case~(C2), \(l\) is nef and there is a prime exceptional divisor \(E\) of class
  \(-l + m\).  The effective cone is generated by \(l\) and \(-l + m\), so some positive
  power of \(L \otimes M\) defines a divisorial contraction \(c \colon X \to Y\) with
  exceptional locus \(E\).

  \textbf{Step~1 (sections of \(L^2\)):} Since \(c_E^\ast L\) restricts to \(H_\Sigma\)
  on \(\Sigma = c(E)\), which is an ample line bundle on the K3 surface \(\Sigma\) (Lemma~7.2
  and Lemma~8.2 of the paper), one has \(H^2(X, M^2) = H^4(X, M^2) = 0\) and
  \(\chi(X, M^2) = P_{RR,X}(0) = 3\) (Lemma~7.3), giving \(h^0(X, M^2) \ge 3\).  The
  reflection that permutes \(l\) and \(m\) is a monodromy operator (Section~3.2); by
  deformation and upper semi-continuity, \(h^0(X, L^2) \ge 2\).

  \textbf{Step~2 (semi-ampleness and Lagrangian fibration):} Since \(\kappa(X, L) > 0\)
  (as \(h^0(X, L^k) \ge 2\) for some \(k\)) and \(\dim X = 4\), the line bundle \(L\) is
  semi-ample by \cite[Theorem~1.5]{fuk}.  For \(k_L\) sufficiently large and divisible,
  the sections of \(L^{k_L}\) define a morphism \(f \colon X \to B\) with connected fibers.
  By \cite{hux}, \(B \cong \mathbb{P}^2\) and \(f\) is a Lagrangian fibration with
  \(f^*\mathcal{O}_{\mathbb{P}^2}(1) \cong L^{k_L}\).
\end{proof}

